% !PRINTED 526
\weberpage{526}
%<<CONT
\gesperrt{it is the symmetric or the alternating group, and}

2. \gesperrt{if it contains two transpositions with a common element, it is the symmetric group.}

For if the cycle $(0,1,2)$ occurs in a group, and the group is moreover transitive, then it contains a permutation $\pi$ by which 0 is carried into 3. In the three-term cycle $\pi^{-1}(0,1,2)\pi$ the figure 3 therefore certainly occurs; it is thus, say, $(0,1,3)$, and accordingly we have in the group
\[
\begin{gathered}
(0,1,2), \quad (0,1,2)^2 = (0,2,1) \\
(0,1,3), \quad (0,1,3)^2 = (0,3,1) \\
(0,1,2)(0,3,1) = (1,2,3) \\
(1,2,3)^2 = (1,3,2) \\
(0,3,1)(0,1,2) = (0,3,2) \\
(0,3,2)^2 = (0,2,3),
\end{gathered}
\]
that is, all the three-term cycles that exist at all. Accordingly all permutations of the alternating group occur in the group in question (§. 153), and it is therefore either exhausted therewith, or it contains in addition a permutation of the second kind and is then the symmetric group.

But if a transitive group contains the two transpositions $(0,1)$, $(0,2)$, then it also contains
\[
(0,1)(0,2) = (0,1,2),
\]
and with it the whole alternating group, and since it also contains a transposition, it is the symmetric group.

If now, then, a group different from $P$ and $Q$ contains first a transposition $(0,2)$, then, if it is transitive, it must also contain $(1,3)$, and accordingly the whole intransitive group of the fourth degree:
\weq{3}{1, \quad (0,2), \quad (1,3), \quad (0,2)(1,3).}
If, then, the group is to be transitive, further permutations must be added. These can be sought only among the four-term cycles or among the pairs of transpositions. But a four-term cycle in which 0 and 2 adjoin one another cyclically cannot occur, because, combined with $(0,2)$, it would give a three-term cycle, such as:
\[
(0,2)(0,2,3,1) = (0,3,1).
\]
